\chapter{Sequências: um primeiro curso}\label{ch:g11:seq}

Uma \omterm{def:g11:seq:sequence}{sequência} é uma lista de números produzida por uma regra: os saldos
sucessivos de uma poupança, o tamanho de uma população ano após ano. Este
capítulo estuda as duas famílias que dominam as aplicações --- as
progressões \emph{aritméticas}, que crescem por passos iguais, e as
progressões \emph{geométricas}, que crescem por razões iguais. A teoria
rigorosa dos limites é desenvolvida no \cref{ch:g12:seq}.

\section{Definir uma sequência}

\begin{definition}[Sequência]\label{def:g11:seq:sequence}
Uma \emph{sequência}\index{sequência} $(u_n)$ associa a cada inteiro
$n \geq 0$ (ou $n \geq 1$) um \omterm{def:g10:numbers:sets}{número real} $u_n$, seu \emph{termo de
índice $n$}. Uma sequência pode ser dada
\begin{itemize}
  \item de forma \emph{explícita}, por uma fórmula de $u_n$ em \omterm{def:g11:func:function}{função} de
        $n$: por exemplo, $u_n = n^2 + 1$;
  \item de forma \emph{recursiva}, pelo primeiro termo e por uma regra
        para passar de cada termo ao seguinte: por exemplo, $u_0 = 3$ e
        $u_{n+1} = 2u_n - 1$.
\end{itemize}
\end{definition}

\begin{example}
Para $u_n = n^2 + 1$: $u_0 = 1$, $u_1 = 2$, $u_2 = 5$ e $u_{10} = 101$
diretamente. Para $u_0 = 3$, $u_{n+1} = 2u_n - 1$: $u_1 = 5$, $u_2 = 9$,
$u_3 = 17$ --- cada termo precisa do anterior; chegar a $u_{10}$ exige
dez passos (ou uma fórmula geral, veja o \cref{exo:g11:seq:11}).
\end{example}

\section{Progressões aritméticas}

\begin{definition}[Progressão aritmética]\label{def:g11:seq:arithmetic}
Uma \omterm{def:g11:seq:sequence}{sequência} é uma \emph{progressão aritmética}\index{progressão
aritmética} de \emph{razão}\index{razão} $d$ se cada termo é obtido do
anterior somando-se $d$:
\[
u_{n+1} = u_n + d \quad \text{para todo } n.
\]
Equivalentemente: a diferença $u_{n+1} - u_n$ é constante, igual a $d$.
\end{definition}

\begin{theorem}[Termo geral]\label{thm:g11:seq:arithgeneral}
Se $(u_n)$ é uma \omterm{def:g11:seq:arithmetic}{progressão aritmética} de primeiro termo $u_0$ e razão
$d$, então
\[
u_n = u_0 + n\,d \quad \text{para todo } n \geq 0,
\qquad\text{e, mais geralmente, } u_n = u_p + (n - p)\,d .
\]
\end{theorem}

\begin{proof}
Para ir de $u_0$ a $u_n$, a regra ``some $d$'' é aplicada $n$ vezes: um
passo dá $u_1 = u_0 + d$, dois passos dão $u_2 = u_0 + 2d$ e, após $n$
passos, cada aplicação contribuiu com um $d$, de modo que
$u_n = u_0 + nd$. (Esse ``e assim por diante'' é tornado rigoroso por
indução no \cref{ch:g12:seq}.) A fórmula geral segue contando os $n - p$
passos de $u_p$ até $u_n$.
\end{proof}

\begin{theorem}[Soma de inteiros consecutivos]\label{thm:g11:seq:intsum}
Para todo inteiro $n \geq 1$:
\[
1 + 2 + \dots + n = \frac{n(n+1)}{2}.
\]
Mais geralmente, uma soma de termos consecutivos de uma \omterm{def:g11:seq:arithmetic}{progressão
aritmética} vale
\[
(\text{número de termos}) \times
\frac{\text{primeiro termo} + \text{último termo}}{2}.
\]
\end{theorem}

\begin{proof}
Escreva a soma $S$ duas vezes, a segunda em ordem inversa, e some coluna
a coluna:
\[
\begin{array}{ccccccccc}
S & = & 1 & + & 2 & + & \dots & + & n\\
S & = & n & + & (n-1) & + & \dots & + & 1\\
\hline
2S & = & (n+1) & + & (n+1) & + & \dots & + & (n+1)
\end{array}
\]
Há $n$ colunas, cada uma somando $n + 1$, logo $2S = n(n+1)$. Para uma
\omterm{def:g11:seq:arithmetic}{progressão aritmética} qualquer o mesmo emparelhamento funciona: primeiro
$+$ último $=$ segundo $+$ penúltimo $= \dots$, porque avançar um passo
na ponta esquerda ($+d$) é compensado por recuar um passo na ponta
direita ($-d$).
\end{proof}

\begin{example}
$1 + 2 + \dots + 100 = \frac{100 \times 101}{2} = 5050$. A soma dos
ímpares $1 + 3 + \dots + 99$ ($50$ termos) é
$50 \times \frac{1 + 99}{2} = 2500$.
\end{example}

\section{Progressões geométricas}

\begin{definition}[Progressão geométrica]\label{def:g11:seq:geometric}
Uma \omterm{def:g11:seq:sequence}{sequência} é uma \emph{progressão geométrica}\index{progressão
geométrica} de \emph{razão}\index{razão} $q \neq 0$ se cada termo é
obtido do anterior multiplicando-se por $q$:
\[
u_{n+1} = q\,u_n \quad \text{para todo } n.
\]
Equivalentemente, quando nenhum termo se anula: o quociente
$\frac{u_{n+1}}{u_n}$ é constante, igual a $q$.
\end{definition}

\begin{theorem}[Termo geral]\label{thm:g11:seq:geomgeneral}
Se $(u_n)$ é uma \omterm{def:g11:seq:geometric}{progressão geométrica} de primeiro termo $u_0$ e razão
$q$, então
\[
u_n = u_0\, q^n \quad \text{para todo } n \geq 0,
\qquad\text{e, mais geralmente, } u_n = u_p\, q^{\,n-p} .
\]
\end{theorem}

\begin{proof}
Mesma contagem de passos do \cref{thm:g11:seq:arithgeneral}: de $u_0$ a
$u_n$, a regra ``multiplique por $q$'' é aplicada $n$ vezes, contribuindo
com um fator $q^n$.
\end{proof}

\begin{theorem}[Soma geométrica]\label{thm:g11:seq:geomsum}
Para todo real $q \neq 1$ e todo inteiro $n \geq 0$:
\[
1 + q + q^2 + \dots + q^n = \frac{1 - q^{\,n+1}}{1 - q}.
\]
\end{theorem}

\begin{proof}
Seja $S = 1 + q + \dots + q^n$. Multiplique por $q$:
$qS = q + q^2 + \dots + q^{n+1}$. Subtraia:
\[
S - qS = \bigl(1 + q + \dots + q^n\bigr)
- \bigl(q + q^2 + \dots + q^{n+1}\bigr) = 1 - q^{\,n+1},
\]
porque cada termo intermediário aparece uma vez em cada soma e se
cancela. Logo $(1 - q)S = 1 - q^{\,n+1}$ e, dividindo por
$1 - q \neq 0$, obtém-se a fórmula.
\end{proof}

\begin{example}\label{ex:g11:seq:chessboard}
$1 + 2 + 4 + \dots + 2^{10} = \frac{1 - 2^{11}}{1 - 2} = 2^{11} - 1 =
2047$: dobrar grãos de arroz nas casas de um tabuleiro de xadrez esgota
qualquer celeiro muito antes da $64^{\text{a}}$ casa, em que o total é
$2^{64} - 1 \approx 1.8 \times 10^{19}$.
\end{example}

\begin{omfigure}
\begin{tikzpicture}
\begin{axis}[omaxis,
  width=10cm, height=5.6cm,
  xlabel={$n$}, ylabel={$u_n$},
  xmin=-0.5, xmax=10.8, ymin=0, ymax=13,
  xtick={5,10}, ytick={2,6}]
  \addplot[omDef, only marks, mark=*, mark size=1.5pt,
    samples at={0,...,10}] {2 + 0.9*x};
  \addplot[omThm, only marks, mark=*, mark size=1.5pt,
    samples at={0,...,10}] {2*1.2^x};
  \node[omDef, font=\small, anchor=east] at (axis cs:9.7,10.9)
    {aritmética};
  \node[omThm, font=\small, anchor=west] at (axis cs:8.6,7.1)
    {geométrica};
\end{axis}
\end{tikzpicture}

{\small Passos iguais contra razões iguais: uma \omterm{def:g11:seq:arithmetic}{progressão aritmética}
($u_{n+1} = u_n + 0.9$, em azul) segue uma reta; uma \omterm{def:g11:seq:geometric}{progressão
geométrica} ($u_{n+1} = 1.2\,u_n$, em vermelho) segue uma curva
exponencial que acaba ultrapassando-a.}
\end{omfigure}

\begin{method}[Reconhecer o tipo de uma sequência]\label{met:g11:seq:recognize}
Calcule $u_{n+1} - u_n$ e simplifique. Se o resultado for uma constante
$d$, a \omterm{def:g11:seq:sequence}{sequência} é uma \omterm{def:g11:seq:arithmetic}{progressão aritmética}. Caso contrário, calcule
$\frac{u_{n+1}}{u_n}$ (termos não nulos) e simplifique: uma constante $q$
indica \omterm{def:g11:seq:geometric}{progressão geométrica}. Se nenhum dos dois for constante, a
\omterm{def:g11:seq:sequence}{sequência} não é de nenhum dos dois tipos --- nunca conclua apenas a
partir dos primeiros termos.
\end{method}

\begin{example}
Para $u_n = 3 \times 5^n$:
$\frac{u_{n+1}}{u_n} = \frac{3 \times 5^{n+1}}{3 \times 5^n} = 5$ para
todo $n$: \omterm{def:g11:seq:geometric}{progressão geométrica} de razão $5$. Para $u_n = n^2$:
$u_1 - u_0 = 1$, mas $u_2 - u_1 = 3$, e $\frac{u_1}{u_0}$ nem sequer está
definido --- não é aritmética nem geométrica.
\end{example}

\section{Monotonicidade}

\begin{definition}[Sequência monótona]\label{def:g11:seq:monotone}
Uma \omterm{def:g11:seq:sequence}{sequência} $(u_n)$ é \emph{\omterm{def:g11:func:monotone}{crescente}} se $u_{n+1} \geq u_n$ para todo
$n$, e \emph{\omterm{def:g10:functions:variations}{decrescente}} se $u_{n+1} \leq u_n$ para todo $n$.
\end{definition}

\begin{method}[Estudar a monotonicidade]\label{met:g11:seq:monotonicity}
Estude o sinal de $u_{n+1} - u_n$. Para \omterm{def:g11:seq:sequence}{sequências} de termos positivos,
pode-se, em vez disso, comparar $\frac{u_{n+1}}{u_n}$ com $1$.
\end{method}

\begin{example}
Uma \omterm{def:g11:seq:arithmetic}{progressão aritmética} é \omterm{def:g11:func:monotone}{crescente} quando $d \geq 0$
($u_{n+1} - u_n = d$) e \omterm{def:g10:functions:variations}{decrescente} quando $d \leq 0$. Uma \omterm{def:g11:seq:geometric}{progressão
geométrica} com $u_0 > 0$ e $q > 1$ é \omterm{def:g11:func:monotone}{crescente}:
$u_{n+1} - u_n = u_0 q^n (q - 1) > 0$; com $u_0 > 0$ e $0 < q < 1$ ela é
\omterm{def:g10:functions:variations}{decrescente}.
\end{example}

\section{Comportamento a longo prazo, informalmente}

O que acontece com $u_n$ quando $n$ fica muito grande? Para uma
\omterm{def:g11:seq:arithmetic}{progressão aritmética} com $d > 0$, os termos $u_0 + nd$ acabam
ultrapassando qualquer número fixado. Para uma \omterm{def:g11:seq:geometric}{progressão geométrica} com
$0 < q < 1$, os termos $u_0 q^n$ encolhem rumo a $0$: multiplicar
repetidamente por $0.9$, digamos, corrói qualquer valor inicial. E, para
$q > 1$, os termos explodem, como no \cref{ex:g11:seq:chessboard}.

\begin{remark}
Essas afirmações podem ser tornadas perfeitamente precisas --- ``os
termos acabam ficando a qualquer distância dada de $0$'' --- e
demonstradas. Essa é a teoria dos \emph{limites}, o tema de abertura do
\cref{ch:g12:seq}.
\end{remark}

\section{Exercícios}

\begin{exercise}[$\star$]\label{exo:g11:seq:1}
Para cada \omterm{def:g11:seq:sequence}{sequência}, calcule $u_1$, $u_2$, $u_3$:
\[
u_n = \frac{n}{n+1}; \qquad
u_0 = 5,\ u_{n+1} = 3u_n - 2; \qquad
u_n = (-1)^n\,n .
\]
\end{exercise}

\begin{exercise}[$\star$]\label{exo:g11:seq:2}
$(u_n)$ é uma \omterm{def:g11:seq:arithmetic}{progressão aritmética} com $u_0 = 7$ e $d = -3$. Calcule
$u_{10}$ e $u_{25}$. $(v_n)$ é uma \omterm{def:g11:seq:arithmetic}{progressão aritmética} com $v_3 = 11$ e
$v_8 = 26$. Encontre a razão e $v_0$.
\end{exercise}

\begin{exercise}[$\star$]\label{exo:g11:seq:3}
$(u_n)$ é uma \omterm{def:g11:seq:geometric}{progressão geométrica} com $u_0 = 5$ e $q = 2$. Calcule
$u_8$. $(v_n)$ é uma \omterm{def:g11:seq:geometric}{progressão geométrica} de termos positivos, com
$v_2 = 12$ e $v_4 = 48$. Encontre a razão e $v_0$.
\end{exercise}

\begin{exercise}[$\star$]\label{exo:g11:seq:4}
Calcule
\[
1 + 2 + 3 + \dots + 500, \qquad
4 + 7 + 10 + \dots + 61, \qquad
1 + \frac12 + \frac14 + \dots + \frac{1}{2^{10}} .
\]
\end{exercise}

\begin{exercise}[$\star$]\label{exo:g11:seq:5}
Determine se cada \omterm{def:g11:seq:sequence}{sequência} é uma \omterm{def:g11:seq:arithmetic}{progressão aritmética}, uma \omterm{def:g11:seq:geometric}{progressão
geométrica} ou nenhuma das duas:
\[
u_n = 4n - 1; \qquad
v_n = \frac{2^n}{3^{n+1}}; \qquad
w_n = n^2 + n .
\]
\end{exercise}

\begin{exercise}[$\star\star$]\label{exo:g11:seq:6}
Um teatro tem $20$ fileiras: $16$ poltronas na primeira fileira, e cada
fileira tem $2$ poltronas a mais que a anterior. Quantas poltronas há na
última fileira? E no teatro inteiro?
\end{exercise}

\begin{exercise}[$\star\star$]\label{exo:g11:seq:7}
Uma população de bactérias dobra a cada hora; ao meio-dia há $500$
bactérias. Quantas há às 20h? Depois de quantas horas inteiras a
população ultrapassa um milhão pela primeira vez? (Resolva testando
potências sucessivas de $2$.)
\end{exercise}

\begin{exercise}[$\star\star$]\label{exo:g11:seq:8}
Todo mês, um poupador deposita $100$ euros em uma conta que paga
$0.2\%$ de juros mensais sobre o saldo existente (os juros são creditados
logo antes do depósito). Seja $c_n$ o saldo imediatamente após o
$n$-ésimo depósito, de modo que $c_1 = 100$ e
$c_{n+1} = 1.002\,c_n + 100$. Calcule $c_2$ e $c_3$ e explique por que
$(c_n)$ não é \omterm{def:g11:seq:arithmetic}{progressão aritmética} nem geométrica.
\end{exercise}

\begin{exercise}[$\star\star$]\label{exo:g11:seq:9}
Estude a monotonicidade das \omterm{def:g11:seq:sequence}{sequências}
\[
u_n = n^2 - 8n \ (n \geq 0), \qquad
v_n = \frac{3^n}{n!}\ (n \geq 1),
\]
em que $n! = 1 \times 2 \times \dots \times n$. (Para $(v_n)$, compare
$\frac{v_{n+1}}{v_n}$ com $1$.)
\end{exercise}

\begin{exercise}[$\star\star$]\label{exo:g11:seq:10}
A soma dos $n$ primeiros termos de uma \omterm{def:g11:seq:arithmetic}{progressão aritmética} com
$u_0 = 3$ e $d = 4$ vale $903$. Encontre $n$. (Monte uma \omterm{def:g10:algebra:equation}{equação} do
segundo grau em $n$ e use o \cref{ch:g11:quad}.)
\end{exercise}

\begin{exercise}[$\star\star\star$]\label{exo:g11:seq:11}
Sejam $u_0 = 3$ e $u_{n+1} = 2u_n - 1$.
\begin{enumerate}
  \item Calcule $u_1, u_2, u_3$ e conjecture uma fórmula para $u_n$.
  \item Seja $v_n = u_n - 1$. Mostre que $(v_n)$ é uma \omterm{def:g11:seq:geometric}{progressão
        geométrica} e dê sua razão e seu primeiro termo.
  \item Deduza uma fórmula explícita para $u_n$ e verifique sua
        conjectura.
\end{enumerate}
\end{exercise}

\section{Problema: A torre de Brama e os coelhos de Fibonacci}

\begin{problem}[{Problema de fim de semana --- duas recorrências
lendárias: a torre que acaba com o mundo, a sequência que cresce
como ouro e o truque auxiliar que domestica empréstimos}]
\label{pb:g11:seq:1}
Duas \omterm{def:g11:seq:sequence}{sequências} reinam no folclore da matemática. Uma conta os
movimentos da torre de Brama --- sessenta e quatro discos de ouro
cuja transferência, diz a lenda, acabará com o mundo. A outra conta
os coelhos de Fibonacci e esconde a razão áurea. Nenhuma das duas
é aritmética, nenhuma é geométrica --- e ambas se rendem às armas
deste capítulo: recorrências, somas geométricas
(\cref{thm:g11:seq:geomsum}) e o truque da \omterm{def:g11:seq:sequence}{sequência} auxiliar do
\cref{exo:g11:seq:11}, que também calcula o seu financiamento.

\medskip
\noindent\textbf{Parte I --- A torre de Brama.}
O quebra-cabeça: $n$ discos de tamanhos \omterm{def:g10:functions:variations}{decrescentes} estão
empilhados no pino A; leve a pilha inteira para o pino C, um disco
por vez, sem nunca colocar um disco maior sobre um menor (o pino B
pode ajudar). Seja $h_n$ o número \omterm{def:g10:functions:extrema}{mínimo} de movimentos.
\begin{enumerate}
  \item Jogue (com moedas) e registre $h_1$, $h_2$, $h_3$.
  \item Explique a estratégia por trás da recorrência
        $h_{n+1} = 2h_n + 1$: o que precisa acontecer antes e
        depois de o maior disco se mover?
  \item Resolva a recorrência com o truque do
        \cref{exo:g11:seq:11}: ponha $v_n = h_n + 1$, mostre que
        $(v_n)$ é uma \omterm{def:g11:seq:geometric}{progressão geométrica} e conclua que
        $h_n = 2^n - 1$.
  \item A torre da lenda tem $64$ discos, e os monges movem um
        disco por segundo. Usando $2^{10} = 1024 \approx 10^3$,
        estime o tempo de transferência em anos (um ano tem cerca
        de $3 \times 10^7$ segundos; compare com o
        \cref{ex:g11:seq:chessboard}, o mesmo gigante em outra
        história). Devemos nos preocupar?
  \item Por que nenhuma estratégia consegue menos que $2^n - 1$
        movimentos? Argumente que \emph{qualquer} solução
        satisfaz $h_{n+1} \geq 2 h_n + 1$: o que tem de ser
        verdade sobre os $n$ discos de cima logo antes e logo
        depois do movimento do disco de baixo?
\end{enumerate}

\medskip
\noindent\textbf{Parte II --- Fibonacci.}
Defina $F_1 = F_2 = 1$ e $F_{n+2} = F_{n+1} + F_n$ (cada termo é a
soma dos dois anteriores --- a regra de contagem de ritmos do
volume do ensino fundamental, agora com seu nome europeu).
\begin{enumerate}[resume]
  \item Liste $F_1$ até $F_{12}$.
  \item Mostre que $(F_n)$ não é \omterm{def:g11:seq:arithmetic}{progressão aritmética} nem
        geométrica, mas que é estritamente \omterm{def:g11:func:monotone}{crescente} a partir de
        $n = 2$ (\cref{met:g11:seq:monotonicity} e a
        recorrência).
  \item Demonstre a identidade das somas
        \[
        F_1 + F_2 + \dots + F_n = F_{n+2} - 1
        \]
        por telescopagem: escreva cada $F_k$ como
        $F_{k+2} - F_{k+1}$ e veja a soma desabar. Verifique-a
        para $n = 6$.
  \item Demonstre a identidade dos quadrados
        $F_1^2 + F_2^2 + \dots + F_n^2 = F_n F_{n+1}$,
        telescopando com
        $F_k F_{k+1} - F_{k-1} F_k = F_k^2$. Verifique para
        $n = 4$. (\omterm{def:g10:functions:function}{Imagem}: quadrados de lados
        $1, 1, 2, 3, 5, \dots$ ladrilham um retângulo --- o
        esqueleto da famosa espiral de Fibonacci.)
  \item A identidade de Cassini afirma que
        $F_{n+1} F_{n-1} - F_n^2 = (-1)^n$. Verifique-a para
        $n = 4, 5, 6$ --- e reconheça o motor do truque do
        quadrado que some, jogado no problema sobre áreas do
        volume do ensino fundamental.
  \item Mostre, a partir da recorrência, que
        $F_{n+2} \geq 2 F_n$: Fibonacci pelo menos dobra a cada
        dois passos --- ela cresce pelo menos tão depressa quanto
        uma \omterm{def:g11:seq:geometric}{progressão geométrica} de razão $\sqrt2$.
  \item Calcule as razões $r_n = \frac{F_{n+1}}{F_n}$ para
        $n = 3$ até $10$ (três casas decimais). Admitindo que
        elas se acomodem em um limite $L$, passe a relação
        $r_{n+1} = 1 + \frac{1}{r_n}$ ao limite e resolva: que
        número do \cref{pb:g10:algebra:1} os coelhos veneram?
\end{enumerate}

\medskip
\noindent\textbf{Parte III --- O truque auxiliar, no banco.}
\begin{enumerate}[resume]
  \item Generalize o \cref{exo:g11:seq:11}: para
        $u_{n+1} = a\,u_n + b$ com $a \neq 1$, ponha
        $\ell = \frac{b}{1 - a}$ (o \omterm{pb:g10:functions:1}{ponto fixo}). Mostre que
        $v_n = u_n - \ell$ é uma \omterm{def:g11:seq:geometric}{progressão geométrica} de razão
        $a$ e conclua que $u_n = a^n (u_0 - \ell) + \ell$.
  \item Um empréstimo: $10\,000$ euros a $1\,\%$ de juros ao mês,
        pagos $300$ euros por mês, de modo que a dívida obedece a
        $d_{n+1} = 1.01\,d_n - 300$. Aplique a questão 13 (o
        \omterm{pb:g10:functions:1}{ponto fixo} primeiro!) para obter uma fórmula explícita
        de $d_n$.
  \item Com uma calculadora, encontre o primeiro mês em que a
        dívida se extingue e o total pago. Quanto custou o
        empréstimo em si?
  \item Uma cidade de $50\,000$ habitantes cresce $2\,\%$ ao ano
        e ainda recebe $1\,000$ recém-chegados:
        $p_{n+1} = 1.02\,p_n + 1000$. Dê a fórmula explícita e a
        população após $10$ anos.
\end{enumerate}

\medskip
\noindent\textbf{Parte IV --- As duas famílias reais.}
\begin{enumerate}[resume]
  \item Calcule $1 + 2 + 3 + \dots + 1000$
        (\cref{thm:g11:seq:intsum} --- a soma do pequeno Gauss,
        do volume do ensino fundamental, agora oficial) e
        $1 + 2 + 4 + \dots + 2^{19}$
        (\cref{thm:g11:seq:geomsum}).
  \item Calcule a soma da \omterm{def:g11:seq:arithmetic}{progressão aritmética}
        $7, 12, 17, \dots, 502$ (quantos termos?).
  \item Plano de poupança: $100$ euros depositados por mês,
        rendendo $0.5\,\%$ ao mês; após o $n$-ésimo depósito o
        saldo é $100\left(1.005^{n-1} + \dots + 1.005 +
        1\right)$. Calcule o saldo após $5$ anos ($n = 60$).
  \item Final --- o kit do domador de \omterm{def:g11:seq:sequence}{sequências}: descrições
        explícitas contra recursivas; as duas famílias reais e
        suas fórmulas de soma; a \omterm{def:g11:seq:sequence}{sequência} auxiliar que
        transforma recorrências afins em geométricas; e
        Fibonacci, primeira cidadã fora das duas famílias,
        domada hoje por identidades e à espera das matrizes (ano
        12) e dos limites para a captura completa. Uma frase para
        cada.
\end{enumerate}
\end{problem}
